\chapter{Capabilities --- Authority Without Ambiguity}

A capability is the answer to one question: \textbf{``May I perform this
operation?''} It is authority, not work. CharlotteOS capabilities are unforgeable,
delegable, and attenuable.

\section{What a capability is}

A capability is a per-address-space handle that names a kernel object and carries
a rights mask. The kernel maintains a \textbf{capability table} for each protection
domain. An entry in this table is the only way a userspace process can reference a
kernel object.

Capabilities have four key properties:

\begin{itemize}
    \item \textbf{Unforgeable.} A process cannot create a capability through
        computation. Capabilities enter a domain only through the bootstrap
        contract or through IPC delegation.
    \item \textbf{Scoped.} A capability is valid only in the domain that holds it.
        Capability index 3 in domain A refers to a different object (or nothing)
        than capability index 3 in domain B.
    \item \textbf{Righted.} A capability carries a rights mask specifying which
        operations are permitted. An endpoint capability might carry only
        \texttt{SEND | CALL} rights; it does not grant \texttt{RECEIVE}.
    \item \textbf{Revocable.} When a kernel object is destroyed, all capabilities
        referencing it become invalid. Subsequent use returns an error, not
        undefined behavior.
\end{itemize}

\section{Object types and their rights}

CharlotteOS defines a fixed set of kernel object types, each with its own rights
model. A capability carries both a type tag and a rights mask.

\begin{longtable}{|p{2.8cm}|p{3.5cm}|p{5.9cm}|}
\hline
\textbf{Type} & \textbf{Example rights} & \textbf{What it authorizes} \\
\hline
Endpoint & RECEIVE, MINT\_CONNECTION & Server-side IPC receive queue \\
Connection & SEND, CALL, DELEGATE & Client-side authority to invoke a service \\
ReplyToken & COMPLETE, DELEGATE & One-shot completion of a pending call \\
CompletionQueue & SUBMIT, WAIT, DRAIN & Submission to / waiting on a CQ \\
MemoryObject & MAP, UNMAP, MOVE, LEND, COPY & Page-backed memory ownership and access \\
Interrupt & ACK, MASK, BIND\_TO\_CQ & Hardware interrupt handling \\
MmioRegion & MAP\_DEVICE & Device register access (MMIO) \\
DmaDomain & PIN, UNPIN & DMA buffer management \\
Device & RESET, CONFIGURE & Device lifecycle management \\
Timer & ARM, CANCEL, BIND\_TO\_CQ & Timer operations \\
\hline
\end{longtable}

\section{Delegation and attenuation}

Capabilities can be \textbf{delegated} to another protection domain through IPC.
When delegating, the sender may \textbf{attenuate} the rights: the receiver gets a
subset of the rights the sender held. A service manager that holds a connection
with \texttt{SEND | CALL | DELEGATE} rights can delegate a client-facing
connection with only \texttt{SEND | CALL}, keeping \texttt{DELEGATE} for itself.

Delegation happens in the kernel, mediated by the IPC path. The sender specifies
the target domain, the capability to delegate, and the reduced rights mask. The
kernel validates that the sender actually holds the capability with at least the
rights being transferred.

\section{Generation-safe lookup}

Capability tables use generation counters to prevent use-after-free through stale
indices. When a capability is revoked, the table slot is freed and its generation
is incremented. Any attempt to use an old index with a stale generation returns
\texttt{InvalidCapability} rather than accessing a different object that happened
to be allocated at the same index.

\section{Bootstrap --- how capabilities enter a domain}

A newly created protection domain starts with an empty capability table. It
receives its initial capabilities through the \textbf{config-page contract}: a
mapped page that the supervisor populates with typed entries before the domain
begins execution. The bootstrap contract delivers at minimum:

\begin{itemize}
    \item A connection to the name service (for service lookup).
    \item For driver domains: device capabilities (MmioRegion, Interrupt).
    \item For service domains: an endpoint on which to listen.
    \item A default CQ handle for submission and waiting.
\end{itemize}

After bootstrap, additional capabilities arrive through IPC: the name service
delegates service connections; the supervisor delegates device grants; other
services delegate whatever the protocol authorizes.

\section{Why capabilities matter for critical software}

\begin{enumerate}
    \item Least privilege is architectural, not advisory: A process
        starts with nothing. Every capability it holds was explicitly granted by
        a trusted component (the supervisor, the name service, or another service
        under a protocol). There is no ambient authority to strip away; the
        default is zero.

    \item Kernel-mediated authority is enumerable: A process's capability table
        records the kernel objects it may use. This makes authority easier to
        audit, but does not by itself account for device side effects, protocol
        state, or implementation defects.

    \item Delegation is mediated: A process cannot grant authority it
        does not hold; it cannot grant rights beyond what it possesses; and every
        delegation crosses the kernel, which enforces both constraints.

    \item Teardown revokes the destroyed domain's capabilities and closes its
        endpoint relationships. Existing tests cover stale connections and
        pending-call reconciliation; this is not a claim of system-wide atomic
        revocation for every future object type or external device effect.

    \item Capability tables are bounded: A process cannot create
        unlimited capabilities to exhaust kernel memory. Submission backpressure
        (\texttt{WouldBlock}) applies when a table is full, forcing the
        application to drain completions before submitting more work.
\end{enumerate}

Capabilities give endpoints and memory objects explicit authority instead of
making them ambient resources addressed only by an index or address.

\endinput
